% Figure: halo slots of one leaf block, their donors and weights (Chapter 2).
% Slots, donors and weights are mesh.halo_entries / halo_donors /
% halo_weights of the generate.jl fixture (generated/halo.tex).
\begin{figure}
  \centering
  \begin{tikzpicture}[x=1.42cm,y=1.42cm]
    \def\fxhaloblock#1#2#3#4#5#6#7#8{}%
    \def\fxslot#1#2#3#4#5#6#7#8{}%
    \def\fxdonor#1#2#3#4#5#6#7#8{}%
    % Pass 1: leaf cells (and the periodic image of root (0,0,0)), clipped.
    \begin{scope}
      \clip (8.45,-1.45) rectangle (14.15,3.55);
      \fill[black!7] (8.45,-1.45) rectangle (14.15,0);
      \def\fxhalocontext#1#2#3#4#5#6{%
        \ifnum#6=1
          \draw[tr cell,tr level #5,dash dot] (#1,#2) rectangle (#3,#4);
        \else
          \draw[tr block,tr level #5] (#1,#2) rectangle (#3,#4);
        \fi}%
      % Generated by docs/techref/figures/generate.jl; do not edit.
\fxhaloblock{10}{0}{12}{2}{1}{1}{1,5,0}{5}
\fxhalocontext{12}{0}{14}{2}{0}{1}
\fxhalocontext{12}{2}{14}{4}{0}{1}
\fxhalocontext{8}{0}{9}{1}{1}{0}
\fxhalocontext{9}{0}{10}{1}{1}{0}
\fxhalocontext{8}{1}{9}{2}{1}{0}
\fxhalocontext{9}{1}{10}{2}{1}{0}
\fxhalocontext{10}{0}{11}{1}{1}{0}
\fxhalocontext{11}{0}{12}{1}{1}{0}
\fxhalocontext{10}{1}{11}{2}{1}{0}
\fxhalocontext{11}{1}{12}{2}{1}{0}
\fxhalocontext{10}{2}{11}{3}{1}{0}
\fxhalocontext{11}{2}{12}{3}{1}{0}
\fxhalocontext{10}{3}{11}{4}{1}{0}
\fxhalocontext{11}{3}{12}{4}{1}{0}
\fxhalocontext{8.5}{2}{9}{2.5}{2}{0}
\fxhalocontext{8.5}{2.5}{9}{3}{2}{0}
\fxhalocontext{9}{2}{9.5}{2.5}{2}{0}
\fxhalocontext{9.5}{2}{10}{2.5}{2}{0}
\fxhalocontext{9}{2.5}{9.5}{3}{2}{0}
\fxhalocontext{9.5}{2.5}{10}{3}{2}{0}
\fxhalocontext{8.5}{3}{9}{3.5}{2}{0}
\fxhalocontext{9}{3}{9.5}{3.5}{2}{0}
\fxhalocontext{9.5}{3}{10}{3.5}{2}{0}
\fxslot{0}{0}{9}{-1}{10}{0}{physical}{1}
\fxdonor{0}{0}{9}{0}{10}{1}{1}{0}
\fxslot{1}{0}{10}{-1}{11}{0}{physical}{1}
\fxdonor{1}{0}{10}{0}{11}{1}{1}{0}
\fxslot{2}{0}{11}{-1}{12}{0}{physical}{1}
\fxdonor{2}{0}{11}{0}{12}{1}{1}{0}
\fxslot{3}{0}{12}{-1}{13}{0}{physical}{1}
\fxdonor{3}{0}{12}{0}{14}{2}{1}{1}
\fxslot{0}{1}{9}{0}{10}{1}{same}{1}
\fxdonor{0}{1}{9}{0}{10}{1}{1}{0}
\fxslot{3}{1}{12}{0}{13}{1}{coarser}{1}
\fxdonor{3}{1}{12}{0}{14}{2}{1}{1}
\fxslot{0}{2}{9}{1}{10}{2}{same}{1}
\fxdonor{0}{2}{9}{1}{10}{2}{1}{0}
\fxslot{3}{2}{12}{1}{13}{2}{coarser}{1}
\fxdonor{3}{2}{12}{0}{14}{2}{1}{1}
\fxslot{0}{3}{9}{2}{10}{3}{finer}{4}
\fxdonor{0}{3}{9}{2}{9.5}{2.5}{0.25}{0}
\fxdonor{0}{3}{9.5}{2}{10}{2.5}{0.25}{0}
\fxdonor{0}{3}{9}{2.5}{9.5}{3}{0.25}{0}
\fxdonor{0}{3}{9.5}{2.5}{10}{3}{0.25}{0}
\fxslot{1}{3}{10}{2}{11}{3}{same}{1}
\fxdonor{1}{3}{10}{2}{11}{3}{1}{0}
\fxslot{2}{3}{11}{2}{12}{3}{same}{1}
\fxdonor{2}{3}{11}{2}{12}{3}{1}{0}
\fxslot{3}{3}{12}{2}{13}{3}{coarser}{1}
\fxdonor{3}{3}{12}{2}{14}{4}{1}{1}

      \fill[tr wall] (8.45,-0.17) rectangle (14.15,0);
      \draw[line width=0.8pt] (8.45,0) -- (14.15,0);
      \draw[line width=0.6pt,dash dot] (12,0) -- (12,3.55);
    \end{scope}
    % Pass 2: the block, its slots and the donor links (donor -> slot).
    \def\fxhalocontext#1#2#3#4#5#6{}%
    \def\fxhaloblock#1#2#3#4#5#6#7#8{%
      \draw[tr root] (#1,#2) rectangle (#3,#4);
      \node[tr small,fill=white,rounded corners=1pt,align=center]
        at ({(#1+#3)/2},{(#2+#4)/2}) {owned cells of $\blk$\\ key $(#7)$};}%
    \def\fxslot#1#2#3#4#5#6#7#8{%
      \draw[tr slot] (#3+0.04,#4+0.04) rectangle (#5-0.04,#6-0.04);
      \pgfmathsetmacro\fxsx{(#3+#5)/2}\pgfmathsetmacro\fxsy{(#4+#6)/2}%
      \node[tr slot #7] (s-#1-#2) at (\fxsx,\fxsy) {};}%
    \def\fxdonor#1#2#3#4#5#6#7#8{%
      \pgfmathsetmacro\fxdx{(#3+#5)/2}\pgfmathsetmacro\fxdy{(#4+#6)/2}%
      \pgfmathsetmacro\fxdist{veclen(\fxdx-\fxsx,\fxdy-\fxsy)}%
      \ifdim\fxdist pt>0.01pt
        \node[circle,fill=black,inner sep=0pt,minimum size=2.4pt] at (\fxdx,\fxdy) {};
        \draw[tr donor link,shorten >=2.8pt] (\fxdx,\fxdy) -- (s-#1-#2);
      \fi}%
    % Generated by docs/techref/figures/generate.jl; do not edit.
\fxhaloblock{10}{0}{12}{2}{1}{1}{1,5,0}{5}
\fxhalocontext{12}{0}{14}{2}{0}{1}
\fxhalocontext{12}{2}{14}{4}{0}{1}
\fxhalocontext{8}{0}{9}{1}{1}{0}
\fxhalocontext{9}{0}{10}{1}{1}{0}
\fxhalocontext{8}{1}{9}{2}{1}{0}
\fxhalocontext{9}{1}{10}{2}{1}{0}
\fxhalocontext{10}{0}{11}{1}{1}{0}
\fxhalocontext{11}{0}{12}{1}{1}{0}
\fxhalocontext{10}{1}{11}{2}{1}{0}
\fxhalocontext{11}{1}{12}{2}{1}{0}
\fxhalocontext{10}{2}{11}{3}{1}{0}
\fxhalocontext{11}{2}{12}{3}{1}{0}
\fxhalocontext{10}{3}{11}{4}{1}{0}
\fxhalocontext{11}{3}{12}{4}{1}{0}
\fxhalocontext{8.5}{2}{9}{2.5}{2}{0}
\fxhalocontext{8.5}{2.5}{9}{3}{2}{0}
\fxhalocontext{9}{2}{9.5}{2.5}{2}{0}
\fxhalocontext{9.5}{2}{10}{2.5}{2}{0}
\fxhalocontext{9}{2.5}{9.5}{3}{2}{0}
\fxhalocontext{9.5}{2.5}{10}{3}{2}{0}
\fxhalocontext{8.5}{3}{9}{3.5}{2}{0}
\fxhalocontext{9}{3}{9.5}{3.5}{2}{0}
\fxhalocontext{9.5}{3}{10}{3.5}{2}{0}
\fxslot{0}{0}{9}{-1}{10}{0}{physical}{1}
\fxdonor{0}{0}{9}{0}{10}{1}{1}{0}
\fxslot{1}{0}{10}{-1}{11}{0}{physical}{1}
\fxdonor{1}{0}{10}{0}{11}{1}{1}{0}
\fxslot{2}{0}{11}{-1}{12}{0}{physical}{1}
\fxdonor{2}{0}{11}{0}{12}{1}{1}{0}
\fxslot{3}{0}{12}{-1}{13}{0}{physical}{1}
\fxdonor{3}{0}{12}{0}{14}{2}{1}{1}
\fxslot{0}{1}{9}{0}{10}{1}{same}{1}
\fxdonor{0}{1}{9}{0}{10}{1}{1}{0}
\fxslot{3}{1}{12}{0}{13}{1}{coarser}{1}
\fxdonor{3}{1}{12}{0}{14}{2}{1}{1}
\fxslot{0}{2}{9}{1}{10}{2}{same}{1}
\fxdonor{0}{2}{9}{1}{10}{2}{1}{0}
\fxslot{3}{2}{12}{1}{13}{2}{coarser}{1}
\fxdonor{3}{2}{12}{0}{14}{2}{1}{1}
\fxslot{0}{3}{9}{2}{10}{3}{finer}{4}
\fxdonor{0}{3}{9}{2}{9.5}{2.5}{0.25}{0}
\fxdonor{0}{3}{9.5}{2}{10}{2.5}{0.25}{0}
\fxdonor{0}{3}{9}{2.5}{9.5}{3}{0.25}{0}
\fxdonor{0}{3}{9.5}{2.5}{10}{3}{0.25}{0}
\fxslot{1}{3}{10}{2}{11}{3}{same}{1}
\fxdonor{1}{3}{10}{2}{11}{3}{1}{0}
\fxslot{2}{3}{11}{2}{12}{3}{same}{1}
\fxdonor{2}{3}{11}{2}{12}{3}{1}{0}
\fxslot{3}{3}{12}{2}{13}{3}{coarser}{1}
\fxdonor{3}{3}{12}{2}{14}{4}{1}{1}

    % Annotations.
    \node[tr small,align=center,fill=white,inner sep=1pt] at (13.45,3.12)
      {periodic image\\ of root $(0,0,0)$};
    \node[tr small,anchor=west,fill=white,inner sep=1pt] at (12.04,3.42) {seam};
    \node[tr small,fill=white,inner sep=1pt] at (9.0,3.35) {$\lev=2$};
    \node[tr small,fill=white,inner sep=1pt] at (8.95,1.0) {$\lev=1$};
    \node[tr small,fill=white,inner sep=1pt] at (13.55,1.45) {$\lev=0$};
    \node[tr small,anchor=north west,fill=black!7,inner sep=1pt] at (8.5,-0.2) {wall, outside $\domain$};
    % Legend.
    \begin{scope}[shift={(8.6,-1.95)}]
      \node[tr slot same] at (0,0) {};
      \node[tr small,anchor=west] at (0.12,0) {same level: copy, $\haloweight{\hc}{b}=1$};
      \node[tr slot coarser] at (2.95,0) {};
      \node[tr small,anchor=west] at (3.07,0) {coarser donor: injection, $\haloweight{\hc}{b}=1$};
      \node[tr slot finer] at (0,-0.38) {};
      \node[tr small,anchor=west] at (0.12,-0.38) {finer donors: mean, $\haloweight{\hc}{b}=\tfrac14$};
      \node[tr slot physical] at (2.95,-0.38) {};
      \node[tr small,anchor=west] at (3.07,-0.38) {physical: boundary rule};
    \end{scope}
  \end{tikzpicture}
  \caption{Halo of width $\halowidth=1$ around the level-one block with key
    $(1,5,0)$ of \cref{fig:mesh:forest} (dashed: its twelve slots
    $\halocells{\blk}$). An arrow runs from each donor centroid to the slot it
    fills by~\cref{eq:mesh:halo-fill}; a same-level slot coincides with its
    donor. On the right the slots lie in the periodic image of a level-zero
    root and copy its single coarse value; the corner slot at the top left
    averages the four level-two cells that contain its quarter-cell sample
    points; the four slots below the wall are set by the boundary rule
    of~\cref{eq:mesh:halo-physical}, which for \code{ZeroGradient} copies the
    marked donor, the cell that contains the clamped slot centre.}
  \label{fig:mesh:halo}
\end{figure}
